% !TEX TS-program = lualatex+se+debug
% !TEX encoding = UTF-8 Unicode

\documentclass[14pt, twoside, a4paper]{extreport}
% -----------------------------------------------------------------------------------
% ------------------------- USTAWIENIA STRONY--------------------------
% -----------------------------------------------------------------------------------
\usepackage{geometry}
\geometry{
	left=10mm,
	right=10mm,
	top=15mm,
	bottom=15mm,
	marginparwidth=0mm,
	marginparsep=0mm
	}
\newcommand{\wstawplik}[1]{\input{../include/#1}}%skraca kod
\wstawplik{headerofficium}
% elementy zmienne co do headerofficium:
\graphicspath{{../img/}} % katalog obrazów
%\svgpath{{../include/}}
\grechangestaffsize{20}%rozmiar wyniku GregorioTEX
% koniec 
\gresetlastline{ragged}% justowanie do lewej
% debugowanie:
% %\iffalse %*** %\fi %***
% \tracingall przed problematycznym kodem - bardzo dużo informacji
% lualatex+se+debug - dodatkowe opcje: -file-line-error i -halt-on-error 
% \errorcontextlines
%\def\mydebug{1}
\begin{document}
\makeatletter
\grechangedim{interwordspacetext}{\unexpanded{\the\fontdimen2\font plus \the\gluestretch\gre@space@skip@interwordspacenotes minus \the\glueshrink\gre@space@skip@interwordspacenotes}}{fixed}%
\makeatother
\pagestyle{empty}
%\pagestyle{plain}	
\large
\begin{spacing}{0.7}
\begin{center}
	\section*{Haec est dies}
	%\small\textit{Fontevraud Gradual\\}
	\small\textit{Cantionale ecclesiasticum (Poznań 1933), p.1}
	\vspace{5mm}
\end{center}	
\end{spacing}
\greannotation{Ant.} 
\greannotation{IV.2} 
%\grecommentary[4mm]{}
\gresetinitiallines{1}\gabcsnippet{(c4) Hæc(ef/ED/ggvFE) est(de) di(efe)es(e) <c><v>*</v></c>(;) quam(e) fe(fe)cit(d) Dó(fg/hvhvGFE)mi(d)nus.(e) (:) 
Hó(eghv)di(h)e(h) Dó(hj/HG/hvGFE)mi(d)nus(e) (;) 
af(fd)flic(f)ti(g)ó(hhv!GF/ede)nem(e) (,) pó(e/fgf)pu(e)li(df) su(d)i(c) res(d)pé(de)xit,(e) (;) et(f) re(d)dem(f)pti(f)ó(gh/hvGFE)nem(d) mi(e)sit.(e) (:) 
Hó(ef/ED/ggvFE)di(d)e(e) mor(eg)tem,(h) (,) quam(h) fé(jh)mi(g)na(gh) ín(g)tu(df)lit(f) (,) fé(fe)mi(d)na(fg/hhvGFE) fu(d)gá(e)vit.(ee) (:) 
Hó(gh!jv)di(j)e(j) De(jvIHj)us(gg) Ho(hg)mo(fe) fac(de)tus,(e) (;) id(e) quod(fvED) fu(fghv)it(hvGFE) per(d)mán(e)sit,(e) (,) et(f) quod(c) non(d) e(efgv)rat(fe) as(d)súm(e)psit.(e) (:) 
Er(g)go(d) ex(g)ór(hgjh)di(g)um(g) (,) nos(jvIG!h)træ(g) re(gvFE)dem(f)pti(d)ó(fg/hvhvGFE/de)nis(e) (;) de(cdev)vó(fvED)te(d) re(gf!gw!hvGF)co(e)lá(de)mus,(e) (:) 
et(g) ex(gf/ghg)sul(fe)té(df)mus(f) di(ef!gvgvFE!f)cén(d)tes:(c) (;) 
Gló(e)ri(d)a(c) ti(d)bi(f) Dó(df!gvgvFE)mi(d)ne.(e) (::) 
<c><i>T.P.</i></c> Al(g)le(gf/gh/hvGFE)lú(de)ia.(e) (::)}
\vspace{5mm}
\setstretch{1}% odstęp między liniami
\small
Ten jest dzień, jenże Bóg uczynił jest. Dzisiaj Bóg na ugabanie luda swego wejrzał jest i posłał
odkupienie. Dzisiaj śmierć, jąż niewiasta zadała, niewiasta oddaliła. Dzisiaj Bóg Człowiek uczynion jest,
jen przed tym był, ostał, a co nie było, wsioł. Tego dla zbawienia naszego bąć obyczaj, nabożnymi słowy
chwalmy rzekąc: Chwała Tobie, Gospodnie, Alleluja!\footnote{Przekład polski datowany na XIV wiek, cyt. za: J. Zathey, Fragment nieznanego modlitewnika staropolskiego z XIV wieku. „Biuletyn Biblioteki Jagiellońskiej” 1963, nr 1/2, s. 37}\newline\newline
Oto jest dzień, który Pan uczynił. Dziś Pan wejrzał na cierpienie swojego ludu i zesłał odkupienie. Dziś śmierć, którą niewiasta sprowadziła, przez niewiastę została przepędzona. Dzisiaj Bóg, który stał się człowiekiem, pozostał Tym, czym był, a przyjął to, czym nie był. Dlatego nabożnie wspominajmy początek naszego odkupienia i radujmy się, mówiąc: Chwała Tobie, Panie. Alleluja, alleluja.\footnote{Tłumaczenie: o. Grzegorz Brodacki OCist}
\end{document}